\documentclass[11pt]{article}

\usepackage[margin=1in]{geometry}
\usepackage[T1]{fontenc}
\usepackage[utf8]{inputenc}
\usepackage{lmodern}
\usepackage{microtype}
\usepackage{amsmath,amssymb,amsthm,mathtools}
\usepackage{enumitem}
\usepackage{booktabs}
\usepackage{tabularx}
\usepackage{xurl}
\usepackage[hidelinks]{hyperref}

\newtheorem{definition}{Definition}
\newtheorem{remark}{Remark}

\newcommand{\TV}{\mathrm{TV}}
\newcommand{\KL}{\mathrm{KL}}
\newcommand{\E}{\mathbb{E}}
\newcommand{\Prb}{\mathbb{P}}

\title{Notation and Minimal Model for the One-True-Archive:\\
\large Keys, state, transcript projections, and the privacy budgets we export}
\author{}
\date{Unified archive note v0.21 (2026-02-28; archive v150)}

\begin{document}
\maketitle

\begin{abstract}
This synthesis note standardizes notation across the One-True-Archive.
It fixes a minimal model (keys, state, transcript, and projections) and names the budget metrics used throughout the series.
The goal is to reduce redundancy: papers may cite this note rather than re-introducing the same random variables, channel semantics, and metric conventions.
\end{abstract}

\paragraph{Scope.}
A compact, series-wide notation sheet and a minimal probabilistic model for Anonymous-DHT lookups.

\paragraph{Imports.}
We treat the backbone mathematics posts as canonical references for mechanism design and lower bounds.
Canonical wiki links (preferred citation targets for the mathematical machinery):
\begin{itemize}[leftmargin=1.2em]
\item \url{[[2026.01.22 - Mathematics: Scheduling and Compiler Techniques for Metadata-Hiding Overlay Lookups]]}
\item \url{[[2026.01.23 - Mathematics: Spectral Anonymity and Optimal Distinguishing Bounds for Random-Walk Delegation in (Possibly Directed) Overlays]]}
\item \url{[[2026.01.25 - Mathematics: Prefix Capacities and Separation Cuts for Shared-Kernel Termination-Time Hiding]]}
\item \url{[[2026.01.26 - Mathematics: Stop-Time Padding Addendum, Max-Mixture Separations and Tail-Sign Deadline Normal Forms]]}
\item \url{[[2026.01.29 - Mathematics: Optimal Poset-Feasible Padding: Knapsack, Transport, and Dual Certificates for TV Privacy]]}
\end{itemize}

\noindent Receipt envelopes (ABOM/OINL) are specified in Anonymity~C; receipt line items in Synthesis~12; digest-bound plan/state adjuncts in Synthesis~18; contract objects in Synthesis~0.\cite{series:abomoinl,series:receiptitems,series:planstateequiv,series:contractobjects}

\noindent For a compact crosswalk from these backbone notes to the receipt- and certificate-facing interfaces in this archive, see Synthesis~24.\cite{series:mathcrosswalk}

\paragraph{Units and logs.}
Unless explicitly marked otherwise, the archive measures privacy and information budgets in \emph{bits}.
Accordingly we write $\log_2$ and $2^{\epsilon}$ for odds-inflation / pointwise-maximal-leakage style caps.
When importing a result stated in nats (i.e.\ using $\ln$ and $e^{\epsilon}$), convert by dividing by $\ln 2$ (so $\epsilon_{\mathrm{bits}}=\epsilon_{\mathrm{nats}}/\ln 2$).
We allow papers to use $\ln$ for convenience in analytic bounds, but any \emph{published budget} should be rendered in bits at the endpoint.

\paragraph{Exports.}
Exports (i) canonical random-variable names ($K,S,T,X$), (ii) the ``projection-first'' transcript interface, and (iii) naming conventions for the budget metrics used in other papers.

\section{Symbol table (quick reference)}
\begin{table}[t]
\centering
\begin{tabularx}{\linewidth}{@{}lX@{}}
\toprule
Symbol & Meaning \\
\midrule
$K$ & Hidden index (DHT key / CID bucket / destination label) the adversary tries to infer. \\
$C=g(K)$ & Coarsening of $K$ (``profile'', ``topic'', committee label) used when many keys share one budget. \\
$S$ & System state (routing tables, caches, congestion, retry counters, persistent randomness). \\
$T$ & Full per-lookup transcript (messages, timing, contacts, retries, receipts). \\
$X=f(T)$ & Declared projection of the transcript that an adversary can reliably aggregate (contact sets, stop-time buckets, witness vectors). \\
$P_{k,s}$ & Conditional law of $X$ (or $T$) given $K=k$ and state $S=s$. \\
$\pi$ & Prior on $K$ (or $C$) used for odds-inflation and identification sizing. \\
\bottomrule
\end{tabularx}
\caption{Canonical random variables and channel notation. Papers may cite this table and only introduce additional symbols.}
\label{tab:symbols}
\end{table}

\paragraph{Artifact policy.}
This archive is source-only; supporting artifacts are available on request as a bundle committed by ABOM/OINL (Anonymity~C) and governed by the artifact policy (Synthesis~6).\cite{series:abomoinl,series:artifactpolicy}
To avoid boilerplate drift, other papers cite this paragraph rather than restating the disclosure policy.

\section{Identifiers and digest conventions (series-wide)}
This series uses two kinds of identifiers: human-stable \emph{names} (for usability) and digest-stable \emph{spec IDs} (for replay and drift control).
We keep the conventions here so other papers do not need to re-state them.

\begin{definition}[Digest ID]
A \emph{digest ID} is a string of the form \texttt{sha256:<hex>}, denoting the SHA-256 digest of a \emph{canonical} serialization of the referenced object (e.g., a JSON plan spec or state declaration).
\end{definition}

\begin{definition}[Receipt envelope IDs]
A published receipt carries top-level IDs:
\texttt{claim\_id} (what claim this is) and \texttt{tw\_id} (which threat-window declaration governs repetition).
Optionally it may carry \texttt{state\_decl\_id} when a claim depends on declared state/measurement surfaces (monitored settings / Regime~III).
Each may be a digest ID.
\end{definition}

\begin{definition}[Replay hook IDs]
A receipt line item may include a replay hook:
\texttt{plan\_id} (a human-stable mnemonic, e.g.\ \texttt{eval-primary-v1}) and \texttt{plan\_spec\_id} (a digest ID binding the referenced plan specification).
Equivalence/drift rules for plan specs and state declarations are centralized in Synthesis~18.
\end{definition}

\begin{definition}[Interface IDs used by line items]
Receipt line items refer to interface families by digest-bound IDs.
In particular \texttt{exposure\_nf\_id} identifies the exposure normal-form (projection schema + conditioning), while \texttt{cond\_iface\_id} may identify a conditional/per-step certificate interface (see Synthesis~21).
Both are digest IDs and should be treated as the ``real'' interface keys for diffing across releases.\cite{series:receiptitems}
\end{definition}

\begin{definition}[Evidence IDs]
When a bound depends on empirical measurement, a line item or certificate may include an \texttt{evidence\_id} that digest-binds a minimal evidence table (counts, confidence parameters, and tool versioning) sufficient for replay.
The semantics of what must appear in such evidence tables are owned by the module exporting the claim (e.g., Anonymity~B for profiling/equalization audits, or Synthesis~21 for conditional per-step budgets).\cite{series:condcert}
\end{definition}


\begin{remark}[Non-redundancy]
If a paper needs only \emph{naming} conventions, cite this note.
If it needs \emph{equivalence} or \emph{recertification} semantics, cite Synthesis~18 (and, when needed, Certified~C for disagreement-gated recertification).
\end{remark}
\section{Minimal model: keys, state, transcript}
\paragraph{Hidden index (key / destination).}
$K\in\mathcal K$ denotes the hidden index the adversary wants to infer.
Depending on context, $K$ may be a concrete DHT key, a CID bucket, a destination class, or a coarse-grained ``topic'' label.
In profiling settings, we also use $C=g(K)$ for an operator-chosen coarsening of keys into ``profiles'' or classes.

\paragraph{System state.}
$S$ is the internal state that affects the distribution of observable behavior:
peer tables, caches, routing scores, congestion state, retry counters, and any persistent randomness.
State may evolve across lookups.
Papers in the AnonDHT State series treat $S$ explicitly and export contracts for controlling its leakage.\cite{series:state1,series:state3}

\paragraph{Transcript.}
A single lookup induces a full (potentially high-dimensional) transcript $T$ capturing messages, timing, peer contacts, abort/retry decisions, and any published receipts.
We avoid committing to a single transcript representation.
Instead, the archive is \emph{projection-first}:
we reason about declared projections $X=f(T)$ that correspond to what an adversary can aggregate reliably.

\paragraph{Projections.}
A projection $X$ may be discrete (contact sets, bucketed times) or continuous (RTT witness vectors).
The trace-interface note (Synthesis~3) provides a family of standard projection schemas; individual papers ``own'' additional projections (e.g., MUCC contact footprints, Congestion-EQ witnesses).\cite{series:traceifaces,series:state2,series:congeq}

\section{Channels and contracts}
A deployment induces a conditional law of transcripts given the hidden key and state.
We write this abstractly as a channel
\[
(K,S) \mapsto T \mapsto X=f(T).
\]
A paper typically exports a \emph{contract} on the $X$-channel of the form:
\begin{itemize}[leftmargin=1.2em]
\item a declared projection schema $X$ (what is observed / published),
\item a budget metric on the family of conditional laws $\{P_{k,s}\}$,
\item an accountant for repetition (horizon-$n$ use, retries, resets, or time windows).
\end{itemize}

\begin{remark}[Non-redundancy principle]
If a paper introduces a new witness tier, it should export a checkable $X$-contract and cite other papers for how to convert that contract into endpoint bounds.
In this archive, odds inflation is ``owned'' by Anonymity~A and the endpoint bridge by Synthesis~5.\cite{series:oddsinflation,series:endpointbridge}
\end{remark}

\section{Budget metrics (naming conventions)}
Different papers use different metrics depending on what is easiest to certify.
Table~\ref{tab:metrics} is the series naming convention.

\begin{table}[t]
\centering
\begin{tabularx}{\linewidth}{@{}lX@{}}
\toprule
Metric name & Meaning / where it appears \\
\midrule
$\TV(P,Q)$ & Total variation distance; common for bounded discrete projections and for TV dual certificates (2026-01-29; Congestion-EQ).\\
$D_\infty^{\delta}(P\Vert Q)$ & Approximate max-divergence: $P(E)\le 2^{\eta}Q(E)+\delta$; used as an equalization/contract surface (Anonymity~B).\\
\emph{Odds inflation} & Multiplicative posterior-odds bound, exported as a single-number budget (Anonymity~A).\\
Endpoint bridge & Conversions among realized odds inflation, PML/MaxL-style quantities, and identification sizing (Synthesis~5).\\
Receipt redaction bound & Leakage from published receipts/certificates as a contraction module (Synthesis~1).\\
\bottomrule
\end{tabularx}
\caption{Series naming conventions for budget metrics. Papers cite this table to avoid re-introducing the same metric taxonomy.}
\label{tab:metrics}
\end{table}

\section{Repetition, retries, and time windows}
\paragraph{Horizon composition.}
If a per-lookup contract yields a budget parameter $b$, then repeated use typically composes to a horizon-$n$ budget $\mathsf{Comp}(b;n)$.
Anonymity~A and Synthesis~5 provide the ``endpoint-friendly'' composition framing.

\paragraph{Retries and fallbacks.}
Retries are treated as a first-class leakage surface.
Operational~A exports a schedule-only wrapper pattern that pins retry behavior and makes its leakage budgetable.\cite{series:operationalA}

\paragraph{Threat windows.}
Several papers use a time-window model: budgets hold for a declared window (24h, epoch length), after which state resets or public randomness refreshes.
Threat-window conventions and window algebra live in Synthesis~4.\cite{series:threatwindows}

\section{How to use this note}
When writing or reading a paper in this archive:
\begin{enumerate}[leftmargin=1.2em]
\item Identify $K$ (what is being hidden) and the adversary-visible projection $X$.
\item Locate (or define) the contract surface: which metric and which repetition accountant.
\item Use the endpoint bridge (Synthesis~5) to translate contract budgets into endpoint anonymity / identification scale.
\item If your claim depends on data-dependent bucket maps or coarsening partitions, apply the selection discipline note (Synthesis~29).\cite{series:selectiondiscipline}
\item If publishing any objects (receipts, certificates), apply receipt-redaction bounds (Synthesis~1) and the release-property ledger discipline (Release~A).\cite{series:receiptreduction,series:releaseproperty}
\end{enumerate}

\begin{thebibliography}{99}\small

\bibitem{series:artifactpolicy}
Anonymous.
\newblock Artifact and Disclosure Policy for the One-True-Archive (Source-Only Public Release).
\newblock Series synthesis note (Synthesis~6), 2026. (This archive.)

\bibitem{series:abomoinl}
Anonymous.
\newblock ABOM and OINL for Anonymity Claims.
\newblock Anonymity series Paper~C, 2026. (This archive.)

\bibitem{series:oddsinflation}
Anonymous.
\newblock Odds Inflation, Privacy Loss, and Endpoint Anonymity Bounds.
\newblock Anonymity series Paper~A, 2026. (This archive.)

\bibitem{series:endpointbridge}
Anonymous.
\newblock Endpoint Metrics Bridge: Odds Inflation, PML/MaxL, and Identification Sizing.
\newblock Series synthesis note (Synthesis~5), 2026. (This archive.)

\bibitem{series:traceifaces}
Anonymous.
\newblock Trace Interfaces for Anonymous DHT Lookups.
\newblock Series synthesis note (Synthesis~3), 2026. (This archive.)

\bibitem{series:receiptitems}
Anonymous.
\newblock Receipt Line Items for Anonymity Claims: A Minimal Tuple Schema for Published Budgets.
\newblock Series synthesis note (Synthesis~12), 2026. (This archive.)

\bibitem{series:condcert}
Anonymous.
\newblock Conditional Per-Step Budget Certificates.
\newblock Series synthesis note (Synthesis~21), 2026. (This archive.)


\bibitem{series:threatwindows}
Anonymous.
\newblock Threat Windows for Anonymous DHT Deployments.
\newblock Series synthesis note (Synthesis~4), 2026. (This archive.)

\bibitem{series:receiptreduction}
Anonymous.
\newblock Receipt Redaction Bounds for Published Claims.
\newblock Series synthesis note (Synthesis~1), 2026. (This archive.)

\bibitem{series:releaseproperty}
Anonymous.
\newblock Release Property Ledger Receipts.
\newblock Release and Destination series Paper~A, 2026. (This archive.)

\bibitem{series:accountingrulebook}
Anonymous.
\newblock Receipt Accounting and Horizon Composition for Anonymous-DHT Budgets.
\newblock Series synthesis note (Synthesis~19), 2026. (This archive.)

\bibitem{series:operationalA}
Anonymous.
\newblock Schedule-Only Retries in Anonymous DHTs.
\newblock Operational series Paper~A, 2026. (This archive.)

\bibitem{series:state1}
Anonymous.
\newblock State-Dependent Anonymity for Anonymous DHT Lookups.
\newblock AnonDHT state series Paper~1, 2026. (This archive.)

\bibitem{series:state2}
Anonymous.
\newblock MUCC: Committee-Contact Privacy.
\newblock AnonDHT state series Paper~2, 2026. (This archive.)

\bibitem{series:state3}
Anonymous.
\newblock Closed-View Auditing and CPPCs.
\newblock AnonDHT state series Paper~3, 2026. (This archive.)

\bibitem{series:congeq}
Anonymous.
\newblock Congestion-EQ: Congestion Witness Equalization.
\newblock Congestion series Paper~1, 2026. (This archive.)


\bibitem{series:contractobjects}
Anonymous.
\newblock Anonymity Contract Objects for Anonymous DHTs: A Practical Interface for Proof-Carrying Claims.
\newblock Series synthesis note (Synthesis~0), 2026. (This archive.)


\bibitem{series:planstateequiv}
Anonymous.
\newblock Digest-Bound Replay Plans and State Declarations: Equivalence, Minimality, and Change Control.
\newblock Series synthesis note (Synthesis~18), 2026. (This archive.)


\bibitem{series:selectiondiscipline}
Anonymous.
\newblock Selection Discipline for Data-Dependent Buckets and Coarsening Maps.
\newblock Series synthesis note (Synthesis~29), 2026. (This archive.)


\bibitem{series:mathcrosswalk}
Anonymous.
\newblock Mathematics Backbone Crosswalk for Anonymous-DHT Receipts.
\newblock Series synthesis note (Synthesis~24), 2026. (This archive.)

\end{thebibliography}

\end{document}
